\subsection*{一步一步建立直观}

\paragraph{1. 先沿计算方向记下中间量}

从输入计算加权和，再计算激活输出，最后得到平方损失。反向传播需要这些前向中间量，因为每个局部导数都要在当前输入处求值。先在右图核对 $z$、输出和目标，才能解释梯度读数为什么具有当前的符号与大小。



\paragraph{2. 从输出误差追问谁该改变}

损失先告诉输出应该增大还是减小，再通过激活函数传到加权和，最后乘以各输入，得到各权重应承担的变化率。这个连乘是链式法则。若一个参数通往损失的路径不止一条，各条路径的贡献还要相加，而不能任选其中一条。



\paragraph{3. 局部通路可能暂时关闭}

当 ReLU 输入严格小于零时，这条通路的局部导数为零，即使当前损失很大，也不会通过它传回非零梯度。在零点本身，经典导数不存在，程序需选定实现约定。这里展示单个节点的机制；大型网络只是把同样的依赖关系按图重复计算。
